\cotefile{57}
\scanpage{271}
- \mspage{9} -
\margin{$n^\circ\ 227$}

entiers. Le corps $k(x)$ étant une extension finie de $k(x_1,\ldots,x_r)$ le groupe des valeurs dans ce corps est aussi canoniquement isomorphe aux entiers. Donc dans le cas général, on peut aussi définir l'\uline{ordre} d'une fonction sur un diviseur premier $W$ : C'est l'entier associé à cette fonction dans le groupe des valeurs de l'anneau de valuation correspondant à $W$.

Le groupe abélien libre engendré par les diviseurs premiers de $V$ est appelé le groupe des \uline{diviseurs} de $V$. Chaque diviseur peut s'exprimer comme somme finie
\[
D=\sum n_iW_i
\]
où les $n_i$ sont des entiers $\ne 0$, et les $W_i$ sont des diviseurs premiers distincts. Le \uline{support} de $D$, i.e. l'ensemble des points dans les $W_i$, sera noté $|D|$.\margin{\uline{\unsure{$\sim W$}}} Si tous les $n_i$ sont $\geq 0$ on dira que $D$ est \uline{positif}. On voit qu'il y a une relation d'ordre partielle entre diviseurs.

Si $z$ est une fonction sur $V$, on peut lui associer le diviseur $\sum_W(\operatorname{ord}_Wz)W$, noté par $(z)$, la somme étant prise sur tous les diviseurs premiers. On sait que la fonction n'a qu'un nombre fini de zéros et de pôles, donc que cette somme est finie. Deux diviseurs $D_1$ et $D_2$ sont \uline{linéairement équivalents} s'ils diffèrent par le diviseur d'une fonction, i.e. s'il existe une fonction $z$ telle que $D_1=(z)+D_2$. Si $D$ est un diviseur, l'ensemble des diviseurs \uline{positifs} linéairement équivalents à $D$ sera noté par $\mathcal{L}(D)$. Il peut être vide. L'ensemble de tous les diviseurs linéairement équivalent à $D$ sera noté $\mathcal{L}^0(D)$.
\mspage{270}
